\section{Sklyaninブロックと次数シフト}\label{sec:elliptic}

ここからは、前節の相対的な存在定理とは異なり、
個々の有限係数の代数を具体的に扱う。
中心三次元と楕円点曲線を同定した後、
次数シフトに含まれる平行移動を追跡し、
可換側とのrotation付き比較を構成する。

\subsection{中心元と点曲線}

\begin{lemma}\label{lem:center}
 \(G_h=x^3+y^3+z^3-2h^2\sym(xyz)\) は、すべての \(h\) で
 \(S_h\) の非零な中心正則元である。
\end{lemma}
\begin{proof}
 \(D=x^3+y^3+z^3\) とする。関係式から
 \[
 [D,x]=h(zy^2z-yz^2y)=h[zy,yz]=h^2[zy,x^2].
 \]
 また六つの順列を同じ順序へ移すと
 \[
 \sym(xyz)=xyz+\frac h2(-x^3+y^3-z^3)
 \]
 となる。この式を \(x\) と交換させ、再び二次関係式を使えば
 \(2[\sym(xyz),x]=[zy,x^2]\) を得る。
 従って \([G_h,x]=0\)。巡回置換で \(y,z\) についても従う。

 非零性は特殊化で失われない形で確認できる。
 自由代数の次数三部分の線形汎関数 \(\ell\) を
 \[
 \ell(x^3)=\ell(y^3)=\ell(z^3)=1,\quad
 \ell(xyz)=\ell(yzx)=\ell(zxy)=h/2,
 \]
 \[
 \ell(xzy)=\ell(zyx)=\ell(yxz)=-h/2
 \]
 とし、ちょうど二文字が等しい語では零とする。
 三つの二次関係式の左右に一文字を掛けた18個の式を \(\ell\) はすべて殺す。
 例えば \(\ell(z(xy-yx-hz^2))=h/2+h/2-h=0\) である。
 従って \(\ell\) は \((S_h)_3\) に降り、\(\ell(G_h)=3\) となる。
 \((a,b,c)=(1,-1,-h)\) はSklyanin代数の12個の非正則パラメータに属さず、
 \(S_h\) はAS-regular domainである \cite[Introduction]{CGWZ}。
 \(h=0\) は多項式環である。よって \(G_h\) は左右正則である。
\end{proof}

標準Sklyanin記法では \((a,b,c)=(1,-1,-h)\) である。
標準中心元の式 \cite[Equation (2.16)]{Walton}
\[
 g=c(a^3-c^3)x^3+a(b^3-c^3)xyz+b(c^3-a^3)yxz+c(c^3-b^3)y^3
\]
に代入し、\(yxz=xyz-hz^3\) を用いると \(g=-hG_h\) となる。
\(h\ne0\)、\(h^6\ne-27\) では点曲線の公式
\cite[Theorem 2.14]{Walton}から
\[
 E_h=V(x^3+y^3+z^3+h^2xyz)\subset\PP^2
\]
を得る。これは滑らかな楕円曲線であり、
\[
 S_h/(G_h)\simeq B(E_h,L,\sigma_h),\qquad L=\OO_{E_h}(1),
\]
と書ける。ここで \(\sigma_h=t_p\) は平行移動である。
原点をinflection point \(o=[1:-1:0]\) に取ると、
多重線形化した関係式を第一成分 \(o\) で評価した核は
\(p=[1:-1:-h]\) である。従ってこの規約で \(\sigma_h=t_p\) となる。
この中心商と楕円曲線の比較は \cite[Sections 3.5--3.6]{EG} も参照。
\(h=0\) は通常のFermat平面三次曲線の斉次座標環である。
この場合は \(p=o\)、\(\sigma_0=\Id\) とする。

\subsection{rotationに含まれる平行移動を消す}

\begin{lemma}\label{lem:elliptic-conjugacy}
 \(E\) を滑らかな楕円曲線、\(L\) を次数三の直線束、
 \(\sigma=t_p\) とする。
 \(C=B(E,L,\sigma)\) のgraded singularity categoryを \(D^b(E)\) と同一視したとき、
 その次数シフトは可換平面三次曲線のrotation
 \(\rho_0=T_{\OO_E}\circ(-\otimes L)\) と共役である。
\end{lemma}
\begin{proof}
 右加群の規約を
 \[
 \Gamma_*(F)_d=H^0(E,F\otimes L_d),\qquad
 L_d=L\otimes\sigma^*L\otimes\cdots\otimes(\sigma^{d-1})^*L
 \]
 と取る。tailsの次数シフトは \(H(F)=\sigma_*(F\otimes L)\) である。
 Gorenstein parameter零のOrlov windowを比較すると
 \[
 \rho_{\mathrm{nc}}=T_{\OO_E}\circ H=t_{-p}^*\rho_0.
 \tag{2}\label{eq:rho-nc}
 \]
 このwindow公式を確認する。特異圏への商を \(q\) とし
 \(\Psi_i=qR\Gamma_{\ge i}\) と置くと、
 \((1)\Psi_1=\Psi_0H\)、\(q(k)=\Psi_0(\OO_E)[1]\)、
 \[
 \Psi_1F\longrightarrow\Psi_0F
 \longrightarrow\RHom(\OO_E,F)\otimes q(k)
 \]
 がある。楕円曲線のSerre双対から
 \(\Psi_1=\Psi_0T_{\OO_E}^{-1}\) となり、\eqref{eq:rho-nc}を得る。

 原点を \(o\)、\(P_q=\OO_E(q-o)\)、\(A_q=(-\otimes P_q)\) とする。
 spherical twistの核は \(I_\Delta[1]\) なので、
 \[
 K_{T_{P_q}}=(P_q^{-1}\boxtimes P_q)\otimes I_\Delta[1].
 \]
 see-sawを両方向のファイバーへ適用すると
 \[
 (P_q^{-1}\boxtimes P_q)\otimes I_\Delta
 \simeq I_{\{y=x-q\}}.
 \]
 よって \(T_{P_q}=t_q^*T_{\OO_E}\) であり、
 \[
 A_q\rho_0A_q^{-1}=t_q^*\rho_0,\qquad
 \rho_{\mathrm{nc}}=A_{-p}\rho_0A_{-p}^{-1}.
 \tag{3}\label{eq:conjugacy}
 \]
 これらは核とwindowによるdg同値として実現される。
 三次周期性の二つの同型の差は \(\HH^0(D^b(E))^\times=\CC^\times\) のscalarであり、
 共役の自然同型をその三乗根で調整すれば周期性も両立させられる。
 \(\HH^{<0}(D^b(E))=0\) なので、この自然同型の等式に追加の高次ホモトピー障害はない。
 後で用いる有限次数差の作用は、この周期性から構成する。
\end{proof}

\begin{corollary}\label{cor:block}
 \(h^6\ne-27\) のとき、三次周期性と次数シフトを保つdg Morita同値
 \[
 \MF^{\ZZ}(S_h,G_h)
 \simeq
 \MF^{\ZZ}\bigl(\CC[x,y,z],x^3+y^3+z^3+h^2xyz\bigr)
 \]
 が存在する。
\end{corollary}
\begin{proof}
 \cite[Theorem 3.6.2]{EG}、Orlov window、
 \cite[Theorem 5.8]{CCKM}によって両辺を \(D^b(E_h)\) と比較し、
 補題\ref{lem:elliptic-conjugacy}を使う。
 MFと特異圏の比較は、factorizationを商環上の両側無限の完全周期分解へ送る
 dg関手で実現する。中心正則元の商はGorensteinであり、complete resolutionは
 MCM加群の安定圏のdg模型を与える。
 CCKMの完全忠実性をすべてのshiftに適用すると、Hom複体の比較がquasi-isomorphismとなる。
 周囲がAS-regularであることによる高次syzygyのfactorization表示が本質的全射性を与える。
 特異圏からwindowへのdg比較は \cite[Section 5, Theorem 5.8]{CT} による。
 \(h=0\) では恒等の可換比較を用いる。
\end{proof}

\section{中心三次式を加える比較と通常次数への降下}\label{sec:TS}

単独の楕円ブロックの比較から中心変数を加えた模型へ進む。
外部積の完全忠実性と生成性を分けて示し、
最後に多重次数から通常の整数次数への降下を行う。

\subsection{二つの次数ブロックと外部積}

中心変数側の三次式をFermat型に制限する必要はない。
この節では \(T=\CC[u_1,\ldots,u_m]\)、\(m\ge1\)、とし、
\(P\in T_3\) は
\[
 J(P)=\left(\frac{\partial P}{\partial u_1},\ldots,
                 \frac{\partial P}{\partial u_m}\right),\qquad
 \dim_\CC T/J(P)<\infty                                      \tag{4a}\label{eq:central-isolated}
\]
を満たすものとする。これは斉次三次式 \(P\) のアフィン臨界点が
原点だけであるという条件である。
\[
 \Gamma=(\ZZ e_0\oplus\ZZ e_1)/\langle3e_0-3e_1\rangle,
 \qquad w=3e_0=3e_1                                       \tag{4b}\label{eq:central-gamma}
\]
と置く。\(S_d\) の次数を \(de_0\)、すべての \(u_i\) の次数を \(e_1\) とする。
したがって中心変数どうしの混合項を含む \(P\) も次数 \(w\) である。

\begin{lemma}[任意の孤立中心三次式に対する外部積]\label{lem:TS}
 \(S\) を連結非負次数付きnoetherian AS-regular \(\CC\)-代数、
 \(g\in S_3\) を中心正則元とする。
 \eqref{eq:central-isolated}の下で、外部積はpretriangulated hullと
 冪等完備化の後にdg Morita同値
 \[
 \MF^\ZZ(S,g)\otL_\CC\MF^\ZZ(T,P)
 \simeq
 \MF^\Gamma(S\otimes_\CC T,g+P)                             \tag{4c}\label{eq:central-TS}
 \]
 を与える。この同値は二つの因子の次数シフトと、
 両者の三次内部シフトを \([2]\) と同一視する周期性を保つ。
\end{lemma}
\begin{proof}
 \(R=S\otimes_\CC T=S[u_1,\ldots,u_m]\)、\(W=g+P\)、
 \(A=R/(W)\)、\(C=S/(g)\) と置く。
 まず \(W\) は中心正則元である。
 実際、\(u\) に関する単項式順序を固定すると、\(P\ne0\) の最高項の係数は
 非零のscalarであり、任意の非零多項式との積の最高項は消えない。
 以下では中心正則元の商に対するMFとgraded singularity categoryの比較
 \cite[Theorem 5.8]{CCKM}を用いる。
 \(\Gamma\to\ZZ\)、\(e_0,e_1\mapsto1\)、は有限の核を持ち、
 正の整数次数を与えるので、有限次数付き自由分解による同じ比較は
 \(\Gamma\)-次数についても適用できる。

 \paragraph{外部積と完全忠実性.}
 factorizationの規約を
 \[
 Q=(Q_1\xrightarrow{\varphi}Q_0
                  \xrightarrow{\psi}Q_1(w)),\qquad
 \psi\varphi=g\Id,\quad\varphi\psi=g\Id
 \]
 とする。これを \(Q^{i+2}=Q^i(w)\) を満たす曲がった複体に展開すると、
 \(Q^0=Q_0\)、\(Q^1=Q_1(w)\) である。
 二つのfactorizationの外部積は、二つの奇微分をKoszul符号付きで足すことにより定める。
 奇微分の外部和の平方は \(g+P\) である。
 二つの \([2]\) を同じ内部シフトに対応させる関係が
 \(3e_0=3e_1\) であり、得られるfactorizationは有限自由である。

 有限自由対象 \(Q,Q'\) と \(K,K'\) に対して、
 係数を \(S\)-次数と \(T\)-次数に分解すると
 \[
 \Hom_{\MF^\Gamma}(Q\boxtimes K,Q'\boxtimes K')
 \simeq
 \operatorname{Tot}\!\left(
   \Hom_{\MF^\ZZ(S,g)}(Q,Q')\otimes_\CC
   \Hom_{\MF^\ZZ(T,P)}(K,K')\right)                         \tag{4d}\label{eq:central-hom}
 \]
 を得る。具体的には、左辺の次数 \(r\) の斉次写像は、
 \(a+b=r\) を満たす右辺の次数 \((a,b)\) の写像の和に一意に分解する。
 二つの内部次数の差を移す操作は \((3,-3)\) の整数倍だけであり、
 これは二つのcohomological degreeを \((2,-2)\) だけ変える操作に対応する。
 微分と合成の符号も外部和のKoszul符号に一致する。
 各因子のHom複体は、有限個の自由加群の次数シフトと係数代数の非負次数性から
 下に有界である。従って各全cohomological degreeに寄与する項は有限個であり、
 \eqref{eq:central-hom}に完備化は必要ない。
 これが外部積のdg完全忠実性を与える。

 \paragraph{中心変数を零とした商への生成性の帰着.}
 \(\MF^\Gamma(R,W)\) の有限自由対象の奇微分を \(D\) とする。
 \(u_i\) は中心変数なので、係数行列の \(u_i\)-微分は \(R\) 上の導分であり、
 \[
 D\,\frac{\partial D}{\partial u_i}
 +\frac{\partial D}{\partial u_i}\,D
 =\frac{\partial P}{\partial u_i}\Id
 \]
 が成り立つ。内部シフトを付けると、これは
 \(\partial_iP:N\to N(2e_1)\) の零ホモトピーである。
 斉次イデアル \(J(P)\) は \((u_1,\ldots,u_m)\)-primaryなので、
 ある整数 \(d\ge2\) に対してすべての \(u_i^d\) が \(J(P)\) に属する。
 斉次係数でこれらを偏微分の線形結合に書けば、
 \(u_i^d:N\to N(de_1)\) も特異圏で零となる。

 \(N\) を有限生成 \(A\)-加群の有界複体で表し、
 \(u_1^d,\ldots,u_m^d\) に沿うKoszul複体を \(K(u^d;N)\) と書く。
 零写像のconeは元の対象を直和因子に持つため、\(N\) は
 \(K(u^d;N)\) の直和因子である。
 より正確には、一変数ごとのKoszul cone関手は完全関手でretractを保つので、
 各 \(u_i^d\) の零作用から得られるretractを順に合成すればよい。
 このKoszul複体のcohomologyは \((u_1^d,\ldots,u_m^d)\) に殺され、
 \[
 (u_1,\ldots,u_m)^{m(d-1)+1}\subset(u_1^d,\ldots,u_m^d)
 \]
 である。従って各cohomology加群は、有限個の \(C=A/(u_1,\ldots,u_m)\)-加群を
 商とする有限filtrationを持つ。標準切断の三角形も使うと、
 右辺の特異圏は \(C\)-加群の制限とその次数シフトで厚く生成される。

 ここで制限が特異圏へ降りることを確認する。
 \(u_1,\ldots,u_m\) は \(A\) の中心正則列である。
 一つの変数 \(u_i\) について、\(R\) で \(u_i a=Wb\) とする。
 \(u_i=0\) とした商では、\(\overline W=g+P|_{u_i=0}\) は正則である。
 特殊化した \(P\) が非零なら最高項のscalar係数を使い、
 零なら \(g\) の正則性を使えばよい。
 従って \(\overline b=0\)、すなわち \(b=u_i b_1\) であり、
 多項式環 \(R\) で \(u_i\) を消去して \(a=Wb_1\) を得る。
 先行する変数を順に零とした同じ議論が正則列を与える。
 よって \(C\) は有限自由Koszul分解を持つ \(A\)-perfect加群であり、
 制限は \(\Dsg(C)\to\Dsg(A)\) を誘導する。

 \paragraph{制限加群の安定化の明示.}
 Euler公式から
 \[
 q_i=\frac13\frac{\partial P}{\partial u_i},\qquad
 P=\sum_{i=1}^{m}u_iq_i
 \]
 と書ける。\(T\) 上の通常のKoszul分解を
 \[
 K_j=T(-je_1)\otimes_\CC\Lambda^j\CC^m,
 \qquad d_K=\sum_i u_i\iota_i
 \]
 とする。\(\iota_i\) は外積代数の縮約、\(\varepsilon_i\) は対応する基底による左外積である。
 \[
 h_K=\sum_iq_i\varepsilon_i:K_j\longrightarrow K_{j+1}(w)
 \]
 と置くと
 \[
 d_Kh_K+h_Kd_K=P\Id,\qquad h_K^2=0.                         \tag{4e}\label{eq:central-koszul}
 \]
 ここで \(q_i\) の次数は \(2e_1\) であり、
 \(K_{j+1}(w)\) のシフトによりすべての表示写像は内部次数零である。

 有限Koszul複体を偶奇で折り畳み、
 \[
 (K_P)_0=\bigoplus_{j\ge0}K_{2j}(jw),\qquad
 (K_P)_1=\bigoplus_{j\ge0}K_{2j+1}(jw)
 \]
 と置く。\(d_K+h_K\) は
 \[
 (K_P)_1\longrightarrow(K_P)_0
             \longrightarrow(K_P)_1(w)
 \]
 を与え、\eqref{eq:central-koszul}によってこれは \(P\) の有限自由factorizationである。

 \(M\) を \(C\)-加群とし、その特異圏での像を表すfactorizationを
 \(Q=(Q_1\xrightarrow\varphi Q_0\xrightarrow\psi Q_1(w))\) とする。
 有限回のsyzygyとcohomological shiftを取れば
 \(M=\operatorname{coker}\varphi\) としてよい。
 実際、元の加群との差に現れる \(C\)-perfect対象は制限後も \(A\)-perfectである。
 \(\psi\varphi=g\Id\) と \(g\) の正則性から \(\varphi\) は単射であり、
 \[
 0\longrightarrow Q_1\xrightarrow\varphi Q_0\longrightarrow M\longrightarrow0
 \]
 は \(S\)-自由分解である。
 これをhomological degree \(1,0\) に置き、
 \(d_Q=\varphi\)、\(s_Q:Q_0\to Q_1(w)\) を \(\psi\) とする。
 \(s_Q\) は \(Q_1\) 上では零とする。

 \(u_iM=0\) として \(M\) を \(R\)-加群に制限したときの有限自由分解は
 \[
 F_a=\bigoplus_{b+c=a}Q_b\otimes_\CC K_c,
 \qquad d=d_Q\otimes1+(-1)^b\otimes d_K
 \]
 である。\(b\in\{0,1\}\) は \(Q\) 側のhomological degreeである。
 \[
 s=s_Q\otimes1+(-1)^b\otimes h_K:F_a\longrightarrow F_{a+1}(w)
 \]
 と置くと、異なる因子の奇作用の反可換性から
 \[
 ds+sd=W\Id,\qquad s^2=0.                                  \tag{4f}\label{eq:central-shamash}
 \]
 が成り立つ。

 このデータの周期化が実際に \(A\)-自由分解を与えることも確認しておく。
 \(F_a=0\) for \(a<0\) とし、
 \[
 \mathcal F_n=
 \bigoplus_{j\ge0}(F_{n-2j}\otimes_R A)(-jw),\qquad n\ge0
 \]
 と置く。微分は同じ \(j\) 内の \(d\) と、\(j\ge1\) について
 \[
 s:(F_{n-2j}\otimes_R A)(-jw)
 \longrightarrow(F_{n-2j+1}\otimes_R A)(-(j-1)w)
 \]
 の和である。\eqref{eq:central-shamash}により平方は零になる。
 各 \(n\) で和は有限である。

 \(j\) による増大濾過を使うと、第一ページは
 \(H_*(F\otimes_R A)\) の内部・homological shiftの和となる。
 \(W\) が正則で \(WM=0\) なので、
 \[
 H_i(F\otimes_R A)=
 \begin{cases}
 M&i=0,\\
 M(-w)&i=1,\\
 0&i\ne0,1.
 \end{cases}
 \]
 次の微分は、\(j\ge1\) の \(H_0\) を \(j-1\) の \(H_1\) へ送る同型である。
 これを \(ds+sd=W\Id\) から直接見られる。
 \(a\in F_0\) が \(M\) の元を表すなら、\(s(a)\) の \(A\) 上の像はcycleであり、
 \(d s(a)=Wa\) である。
 短完全列
 \[
 0\longrightarrow F(-w)\xrightarrow{W}F
      \longrightarrow F\otimes_R A\longrightarrow0
 \]
 のconnecting mapはそのcycleを元の \(a\) に送る。
 従ってこの微分は、上の \(H_1\simeq M(-w)\) の標準同定の下で恒等である。
 隣接する二項はすべて相殺し、\(j=0\) の \(H_0=M\) だけが残る。
 濾過は各全homological degreeで有限なので収束し、
 \(\mathcal F\) は \(M\) の \(A\)-自由分解となる。

 \(F\) は有限複体なので、この分解の十分先は二周期的である。
 そのfactorizationは
 \[
 \widehat F_0=\bigoplus_{j\ge0}F_{2j}(jw),\qquad
 \widehat F_1=\bigoplus_{j\ge0}F_{2j+1}(jw),
 \qquad d+s
 \]
 による折り畳みであり、Koszul符号と内部シフトを含めて
 \(Q\boxtimes K_P\) に一致する。
 実際、\(\widehat F_0\) は
 \(Q_0\otimes(K_P)_0\oplus Q_1\otimes(K_P)_1(w)\)、
 \(\widehat F_1\) は
 \(Q_1\otimes(K_P)_0\oplus Q_0\otimes(K_P)_1\) である。
 従って制限加群の安定化は
 \[
 \operatorname{stab}_A(M)\simeq Q\boxtimes K_P.               \tag{4g}\label{eq:restriction-mf}
 \]
 ここで全体のshiftにも曖昧さはない。
 十分大きい \(\ell\) について
 \(\mathcal F_{2\ell}=(\widehat F_0\otimes_R A)(-\ell w)\)、
 \(\mathcal F_{2\ell+1}=(\widehat F_1\otimes_R A)(-\ell w)\) であり、
 そのcokernelは \(\Omega_A^{2\ell}M\) である。
 \(\Omega_A^{2\ell}M\simeq M[-2\ell]\) と \((w)\simeq[2]\) によって
 この内部シフトはちょうど打ち消される。

 最後に、\(\Gamma/\ZZ e_0\simeq\ZZ/3\) であるため、
 任意の \(\Gamma\)-次数付き \(C\)-加群は、
 三つの次数類 \(0,e_1,2e_1\) に属する部分に分解する。
 各部分は通常の \(\ZZ\)-次数付き \(C\)-加群の内部シフトであり、
 \eqref{eq:restriction-mf}の右辺の \(T\) 側を \((j)\)、\(j=0,1,2\)、
 だけシフトすればすべて得られる。
 前段のKoszul filtrationと合わせて、外部積の像は右辺全体を厚く生成する。
 これで本質的全射性が従う。
 周期性および二つの次数シフトとの両立は、ここで用いた外部和の構成そのものから従う。
\end{proof}

\begin{remark}[一変数の場合の明示行列]\label{rem:central-one-variable}
 \(m=1\)、\(P=u^3\) では \(K_P=(u,u^2)\) である。
 上の分解とホモトピーは
 \[
 Q_1(-e_1)\xrightarrow{(-u,\varphi)^t}
 Q_1\oplus Q_0(-e_1)\xrightarrow{(\varphi,u)}Q_0,
 \]
 \[
 s_0=(\psi,u^2)^t,\qquad s_1=(-u^2,\psi)
 \]
 となる。折り畳んだfactorizationは
 \[
 \alpha=
 \begin{pmatrix}\varphi&u\\-u^2&\psi\end{pmatrix}:
 Q_1\oplus Q_0(-e_1)\longrightarrow Q_0\oplus Q_1(w-e_1),
 \]
 \[
 \beta=
 \begin{pmatrix}\psi&-u\\u^2&\varphi\end{pmatrix}:
 Q_0\oplus Q_1(w-e_1)\longrightarrow
 Q_1(w)\oplus Q_0(w-e_1).
 \]
 中心性により \(\alpha\beta=\beta\alpha=(g+u^3)\Id\) である。
 これらの式は、一般の \(K_P\) による安定化の符号と次数シフトの具体的な検算になる。
\end{remark}

\subsection{有限群による整数次数の回復}

\(\Gamma\to\ZZ\)、\(e_0,e_1\mapsto1\)、の核は
\(T_\Gamma\simeq\ZZ/3\) であり、\(e_1-e_0\) で生成される。
分裂 \(1\mapsto e_0\) の下で、その双対群
\(H_\Gamma=\Hom(T_\Gamma,\CC^\times)\simeq\mu_3\) は
\(S\) 上恒等、すべての \(u_i\) を同時に三乗根倍する作用である。
\(P\) は斉次三次式なのでこの作用は \(W=g+P\) を保ち、
\[
 \MF^\Gamma(R,W)
 \simeq\bigl(\MF^\ZZ(R,W)\bigr)^{H_\Gamma}.                \tag{5}\label{eq:eq}
\]
有限群の位数は可逆である。
同変な有限次数付き自由加群は、連結な整数次数と有限群表現の半単純性により
同変な斉次基底を持つので、この比較は有限自由factorizationのdg模型で実現する。
さらに双対群 \(T_\Gamma\) で同変化すると、元の通常次数の圏が回復される
\cite[Theorem 4.2, Section 8, Lemma 8.6, Theorem 8.9]{Elagin}。
冪等完備化はこの比較の一部として取る。

外部積 \eqref{eq:central-TS} の左辺で双対作用の生成元は
\[
 \rho_T\rho_S^{-1}
\]
である。ここで \(\rho_S,\rho_T\) はそれぞれの因子の内部次数一のシフトである。
その三乗を恒等と同一視する自然同型は、両因子の
\((3)\simeq[2]\) から得られる。
二つのperiodicityを含めて \(\rho_S\) を保つdg同値はこの作用と両立し、
\eqref{eq:eq}と双対同変化により通常次数へ降りる。
この際 \(e_0\) のシフトも保存されるので、通常次数のdegree shift \((1)\) が保存される。
有限群作用のcoherenceはfactorizationの内部シフトのcoherenceから引き継ぎ、
単に自己同値の三乗が同型になるという存在命題だけから群作用を作るのではない。

\subsection{Sklyaninブロックへの適用}

\begin{proof}[定理\ref{thm:global}の証明]
 補題\ref{lem:TS}を \((S_h,G_h)\) と
 \((\CC[x,y,z],x^3+y^3+z^3+h^2xyz)\) のそれぞれに適用する。
 系\ref{cor:block}は第一因子の次数シフトと三次周期性を保ち、
 \(\MF^\ZZ(T,P)\) 側では恒等を用いる。
 前節の降下によって通常次数のdegree shiftを保つdg Morita同値
 \[
 \MF^\ZZ(S_h[u],G_h+P)
 \simeq
 \MF^\ZZ\!\left(\CC[x,y,z,u],x^3+y^3+z^3+h^2xyz+P\right)
 \]
 を得る。体上のOrlov同値と \cite{CCKM} により、これは残余圏の比較である。

 可換方程式の \(u_i\)-偏微分がすべて零なら、
 \eqref{eq:central-isolated}によって全 \(u_i=0\) である。
 \(h^6\ne-27\) の下では残る平面三次式のアフィン臨界点も原点だけである。
 よって射影超曲面 \(Y_{h,P}\) は滑らかである。
 次元は \(m+1\ge3\) なので、一意性は \cite[Theorem 1.7]{LZ} による。
\end{proof}



\begin{corollary}[平行移動を伴う一般の楕円ブロック]\label{cor:elliptic-central-cubic}
 \(S\) をHilbert series \((1-t)^{-3}\) を持つnoetherianな
 二次三次元AS-regular \(\CC\)-代数、\(g\in S_3\) を中心正則元とする。
 滑らかな平面三次曲線 \(E\)、\(L=\OO_E(1)\)、および
 ある原点に関する平行移動 \(\sigma=t_p\) に対して、次数付き同型
 \[
 S/(g)\simeq B(E,L,\sigma)
 \]
 が存在すると仮定する。
 \(f_E\in\CC[x,y,z]_3\) を \(E\subset\PP^2\) の方程式とし、
 \(P\in\CC[u_1,\ldots,u_m]_3\)、\(m\ge2\)、は
 \eqref{eq:central-isolated}を満たすとする。
 このとき \(S[u]/(g+P)\) のrotation付き残余圏の滑らかな可換モデルは
 \[
 V(f_E(x,y,z)+P(u))\subset\PP^{m+2}
 \]
 であり、射影同型を除いて一意である。
\end{corollary}
\begin{proof}
 Gorenstein parameter零のOrlov同値により
 \(\MF^\ZZ(S,g)\) は \(D^b(E)\) と同値であり、
 次数シフトは補題\ref{lem:elliptic-conjugacy}の
 \(T_{\OO_E}\sigma_*(-\otimes L)\) に対応する。
 同補題の直線束による共役とperiodicityのscalar調整により、
 \(\MF^\ZZ(\CC[x,y,z],f_E)\) との、次数シフトと三次周期性を保つdg同値を得る。
 あとは補題\ref{lem:TS}と有限次数差による降下を適用すればよい。
 滑らかさと一意性は定理\ref{thm:global}と同じ証明による。
 方程式 \(f_E\) の非零scalar倍による違いは、
 分離された \(x,y,z\) の同時三乗根倍によって射影同型に吸収される。
\end{proof}

\begin{remark}[この拡張の範囲]
 定理\ref{thm:global}では中心変数間の任意の三次混合項を許すが、
 非可換ブロックの生成元と中心変数の混合項は含めない。
 別の非可換ブロックについては、その生成元による通常の偏微分が代数の導分とは限らず、
 本節の \(J(P)\)-primaryな支持への帰着は適用できない。
 この点は独立した生成性の問題として残る。
 本節の拡張の先行研究上の新規性は未確認である。
\end{remark}

